\section{Introduction}
\label{sec:intro}

The mass of a main-sequence star is the quantity that most directly sets its
luminosity and evolution, but it can be measured only for stars whose orbits
can be characterised.  Empirical mass--luminosity relations (MLRs) of low-mass
stars are therefore anchored to a few hundred eclipsing and astrometric
binaries with precise dynamical masses
\citep{Delfosse2000,Torres2010,Mann2019}.  These samples are dominated by
stars of near-solar composition.  \citet{Mann2019} found no significant
metallicity dependence of the $M_{K_s}$--mass relation, where models predict
the dependence to be weak; in optical bands such as \gaia\ $G$ the predicted
dependence is much stronger, and it is in practice supplied by
stellar-evolution models such as PARSEC \citep{Bressan2012PARSEC}.  At fixed absolute magnitude, a metal-rich star is
expected to be more massive than a metal-poor one, because higher opacity
lowers the luminosity at fixed mass.  The amplitude of this effect, which
reaches tens of per cent across the metallicity range of the Galactic disc,
has not been measured dynamically for large samples.

Wide binaries offer a statistical route to such a measurement.  With projected
separations of $10^2$--$10^4$~AU, their orbital periods greatly exceed the
\gaia\ baseline, so individual orbits cannot be solved.  The projected
relative velocity $\Delta v_\perp$ and projected separation $r_\perp$ of a
bound pair are, however, related to the total mass by Kepler's third law.
For a population with a known distribution of eccentricities, phases, and
orientations, the combination $u=\Delta v_\perp\sqrt{r_\perp}$ scales as
$\sqrt{M_{\rm tot}}$ times a dimensionless variable whose distribution is fixed
by orbital geometry.  \citet{Hwang2024DynamicalMasses} used this statistic to
derive dynamical masses across the Hertzsprung--Russell diagram from the
catalogue of \citet{ElBadry2021MillionBinaries}.  The same statistic underlies
tests of gravity in the low-acceleration regime with wide binaries
\citep{Chae2023,Banik2024}, which rely on an assumed MLR to predict the
Newtonian velocity scale.

The components of a wide binary formed together and share their birth
composition.  Metallicities of the individual stars can now be inferred for
millions of \gaia\ sources from the low-resolution BP/RP (XP) spectra
\citep{GaiaCollaboration2023DR3,Li2026JCaPS}.  Each binary thus supplies two
noisy estimates of a single metallicity and one dynamical constraint on the sum
of two masses.  Combined over a large sample, these data constrain the MLR as a
function of both luminosity and metallicity.

Turning this into a measurement requires care at three points.  First, the XP
metallicities have magnitude-dependent systematics and heavy-tailed errors,
which must be calibrated before they are used to assign binaries to
metallicities.  Second, the mass enters the likelihood only through the scale
of a broad velocity distribution, so the result depends on the assumed shape of
that distribution and on the treatment of contaminants such as chance
alignments and hierarchical triples.  Third, a flexible MLR must remain
physically ordered: brighter stars must be more massive at fixed composition.

In this paper we present a modular Bayesian framework that addresses these
points and apply it to 14\,876 main-sequence wide binaries within 1~kpc.  A
first module infers a posterior for the shared metallicity of each system from
the component XP metallicities and the PARSEC colour--magnitude relation.  A
second module represents the MLR as a hard-monotone B-spline correction to the
PARSEC mass surface and fits it to the wide-binary velocities while
marginalising over each system's metallicity posterior.  We describe the
sample in Sect.~\ref{sec:data}, the model in Sect.~\ref{sec:methods}, and the
results in Sect.~\ref{sec:results}.  Section~\ref{sec:discussion} discusses
the degeneracy between the mass scale and the velocity distribution and the
limitations of the present analysis, and Sect.~\ref{sec:conclusions}
summarises our conclusions.
